% Cronograma semestral — Meteorologia Prática (I e II)
\documentclass[11pt]{article}
\usepackage{../comum/preambulo}

\title{Cronograma do Curso\\
{\large Meteorologia Prática --- baseado em R.~Stull,
\emph{Practical Meteorology} (CC BY-NC-SA 4.0)}}
\author{}
\date{}

\begin{document}
\maketitle
\thispagestyle{fancy}

\section*{Premissas}

\begin{itemize}
  \item \textbf{Ritmo}: 2 encontros semanais de 2\,h (4\,h/semana);
        cada capítulo ocupa 2 encontros (Aula 1 e Aula 2, conforme o
        plano no topo das notas de cada capítulo).
  \item \textbf{Carga}: 23 capítulos $\Rightarrow$ o curso é uma
        \emph{sequência de dois semestres} de $\sim$15 semanas:
        \textbf{Meteorologia Prática I} (Caps.~1--11: fundamentos até
        a circulação geral) e \textbf{Meteorologia Prática II}
        (Caps.~12--23: sistemas de tempo, camada limite, previsão,
        clima e óptica).
  \item \textbf{Listas}: a lista do capítulo $N$ (T1--T5 escritos $+$
        N1--N4 no notebook) é liberada ao fechar o capítulo e
        entregue \textbf{uma semana depois}
        (\texttt{listas/listas\_exercicios.pdf}, uma página por
        capítulo).
  \item \textbf{Notebooks em aula}: os ``Casos'' indicados entram na
        própria aula (10--20\,min), conforme o guia do professor.
\end{itemize}

\begin{noquadro}[avaliação sugerida (cada semestre)]
2 provas escritas (50\%) --- com consulta ao livro, no espírito do
Stull;\\
listas T $+$ notebooks N (30\%) --- média das entregas semanais;\\
projeto final (20\%) --- estender um exercício N em um estudo próprio
(dados reais, variação de parâmetros, caso brasileiro), entregue como
notebook executado $+$ apresentação de 10\,min na última semana.
\end{noquadro}

\clearpage
\section*{Meteorologia Prática I --- Caps.~1--11 (15 semanas)}

\begin{center}
\begin{tabular}{clll}
\hline
Sem. & Encontros & Conteúdo & Entregas/observações \\
\hline
1 & 1.1 e 1.2 & Cap.~1 Fundamentos & --- \\
2 & 2.1 e 2.2 & Cap.~2 Radiação & Lista 1 \\
3 & 3.1 e 3.2 & Cap.~3 Termodinâmica & Lista 2 \\
4 & 4.1 e 4.2 & Cap.~4 Vapor d'Água & Lista 3 \\
5 & 5.1 e 5.2 & Cap.~5 Estabilidade & Lista 4 \\
6 & revisão / \textbf{P1} & Caps.~1--5 & Lista 5 \\
7 & 6.1 e 6.2 & Cap.~6 Nuvens & --- \\
8 & 7.1 e 7.2 & Cap.~7 Precipitação & Lista 6 \\
9 & 8.1 e 8.2 & Cap.~8 Satélites e Radar & Lista 7 \\
10 & 9.1 e 9.2 & Cap.~9 Boletins e Cartas & Lista 8;
9.2 $=$ \emph{prática de carta impressa} \\
11 & 10.1 e 10.2 & Cap.~10 Forças e Ventos & Lista 9 \\
12 & 11.1 e 11.2 & Cap.~11 Circulação Geral & Lista 10 \\
13 & revisão / \textbf{P2} & Caps.~6--11 & Lista 11 \\
14 & substitutiva / projetos & apresentações finais & projeto final \\
15 & margem & reposição/feriados & --- \\
\hline
\end{tabular}
\end{center}

\paragraph{Sugestões de calendário.} Rodar o Caso~1 do Cap.~14 (CAPE
do dia) \emph{não} pertence a este semestre, mas o Caso~4 do Cap.~1
(sondagem de Campo de Marte) e a prática de carta do Cap.~9 rendem
mais em dias de tempo ativo — vale trocar a ordem dos encontros da
semana se a meteorologia colaborar. A demonstração do tanque girante
(guia do Cap.~11) precisa de montagem prévia.

\clearpage
\section*{Meteorologia Prática II --- Caps.~12--23 (15 semanas)}

\begin{center}
\begin{tabular}{clll}
\hline
Sem. & Encontros & Conteúdo & Entregas/observações \\
\hline
1 & 12.1 e 12.2 & Cap.~12 Massas de Ar e Frentes & --- \\
2 & 13.1 e 13.2 & Cap.~13 Ciclones Extratropicais & Lista 12 \\
3 & 14.1 e 14.2 & Cap.~14 Tempestades & Lista 13 \\
4 & 15.1 e 15.2 & Cap.~15 Perigos de Tempestades & Lista 14 \\
5 & 16.1 e 16.2 & Cap.~16 Ciclones Tropicais & Lista 15 \\
6 & 17.1 e 17.2 & Cap.~17 Ventos Regionais & Lista 16 \\
7 & revisão / \textbf{P1} & Caps.~12--17 & Lista 17 \\
8 & 18.1 e 18.2 & Cap.~18 Camada Limite & --- \\
9 & 19.1 e 19.2 & Cap.~19 Dispersão de Poluentes & Lista 18 \\
10 & 20.1 e 20.2 & Cap.~20 Previsão Numérica & Lista 19 \\
11 & 21.1 e 21.2 & Cap.~21 Clima Natural & Lista 20 \\
12 & 22.1 e 22.2 & Cap.~22 Óptica Atmosférica & Lista 21 \\
13 & 23.1 e 23.2 & Cap.~23 Mudança Climática (extra) & Lista 22 \\
14 & revisão / \textbf{P2} & Caps.~18--23 & Lista 23 \\
15 & substitutiva / projetos & apresentações finais & projeto final \\
\hline
\end{tabular}
\end{center}

\paragraph{Sugestões de calendário.} Agendar os Caps.~12--15 para o
semestre que contém o \emph{inverno/primavera} austral maximiza a
chance de frentes, ciclogêneses do Prata e convecção severa reais
durante as aulas — o meteograma do Cap.~12 e o CAPE do dia (Cap.~14,
Caso~1 com a sondagem das 12\,UTC) devem ser rodados com dados do
próprio dia sempre que possível. O episódio de poluição do Cap.~19
conversa com o inverno estagnado de São Paulo (dados CETESB). A aula
23.2 fecha o curso com a retrospectiva sugerida no guia.

\section*{Trilha compacta (opcional, 1 semestre)}

Para uma disciplina única de tópicos (4\,h/semana, 15 semanas):
1 encontro por capítulo nos Caps.~1--5 e 10--14 (núcleo físico
indispensável, 10 encontros), 2 encontros para o Cap.~20, e os demais
capítulos como seminários dos alunos (2 por encontro, com os
notebooks como roteiro). Provas nas semanas 8 e 15. As notas de aula
suportam o corte: as seções marcadas como Aula 1 carregam o essencial
de cada capítulo.

\vspace{1em}
\noindent\rule{\linewidth}{0.4pt}\\
{\scriptsize Material derivado de R.~Stull, \emph{Practical
Meteorology} (CC BY-NC-SA 4.0); este documento herda a mesma licença.}

\end{document}
